\section{相关工作}

\subsection{元推理与随时规划中的计算与行动权衡}

有限理性与元推理理论旨在回答智能体在面对计算资源约束时应如何分配内部推理时间与外部物理行动\cite{horvitz1987reasoning, russell1991right}。经典单步元推理通过估计额外一次搜索迭代对当前最优动作期望效用的提升量，并在计算期望增益低于时间延迟成本时终止搜索\cite{lin2015metareasoning}。在移动机器人运动规划与并发控制中，由于机体移动与后台计算可以并行执行，研究者提出了边移动边思考的并发控制框架与形式化元推理模型，以权衡继续规划带来的路径质量改善与推迟动作切换造成的跟踪偏差\cite{xiao2020thinking, elboher2023formal}。针对单步解析估计难以捕捉长周期代价的问题，另一类工作利用历史任务轨迹训练停止策略或计算预算选择器，根据环境特征直接预测不同算力预算的收益\cite{sung2021quit, magnino2026metareasoning}。近期在大模型智能体研究中也发现了类似的推理与交互权衡现象，过度拉长内部思考链反而会挤占环境交互机会并降低任务成功率\cite{shen2025thinking, cuadron2025overthinking, hunt2021formalizing}。然而现有元推理方法在评估额外规划收益时，普遍直接信任规划器在自身搜索空间或内部模型上的自评得分，或依赖大量历史任务数据训练预算选择器。本文则指出在不确定海洋场中规划器自评得分受优化器诅咒严重污染，并利用留出想象海洋在无需历史任务的条件下实现了计算价值的无偏评估。

\subsection{海洋自适应采样与机器人信息收集}

自主水下航行器在海洋藻华追踪与锋面观测中的应用，推动了从规则覆盖向信息驱动自适应采样的演进\cite{fossum2019}。传统海洋调查多采用割草机式全覆盖航线以均匀扫测目标海域\cite{galceran2013coverage}，但在大范围赤潮应急巡测中，受航行器巡航速度与24小时作业窗口限制，全覆盖航线仅能扫过窗口的一小部分水域。为提高热区定位效率，信息收集算法通常利用高斯过程回归维护环境场的均值与方差，并采用贪心策略、上置信界采样\cite{srinivas2012gpucb}或汤普森采样\cite{russo2018thompson}引导航行器驶向高浓度或高不确定性位置，或通过采样树搜索规划多步信息收集路径\cite{hollinger2014sampling}。与此同时，针对东海原甲藻赤潮的卫星遥感研究揭示了藻华在晚春至初夏的高发规律及随潮汐往复漂移的日内变化特征\cite{lou2014diurnal, zhao2025distinguishing}，为巡测提供了物理背景。然而现有海洋自适应采样方法均假定机载路径规划的计算开销可忽略或固定不变，未考虑高功耗芯片深想与推进电机共享机载电池及作业时间的耦合约束。本文则将机载算力速度、芯片功耗与原位复测统一纳入随水漂移参考系下的决策空间，系统揭示了多步航线规划与机动采样之间的收益边界。

\subsection{优化器诅咒与卫星数据驱动的集合数据同化}

在随机优化与强化学习中，使用同一批有限样本同时开展策略寻优与价值估计会引入系统性的正向偏差，这一现象在决策分析中被称为“优化器诅咒”\cite{smith2006curse}。当候选方案数量或搜索预算增大时，最大化算子会放大样本估计中的随机正误差，导致选出方案的事后真实表现显著低于优化器的内部预期，而双估计方法通过将动作选择与价值评估解耦至独立样本集来消除这一过估计\cite{vanhasselt2010double}，且模型误差还会进一步限制有效规划视野\cite{jiang2015horizon}。在海洋环境重构方面，业务化水色遥感利用经验正交函数缺失数据插补方法生成逐日无云叶绿素场\cite{beckers2003dineof}，结合多尺度融合海表温度产品\cite{chin2017mur}与水色传感器实测精度评估\cite{ye2021cocts}，以及自监督联合嵌入预测架构等时空模型\cite{assran2025vjepa2}，共同刻画了表层海洋的演变尺度。为融合现场观测，序列集合调整卡尔曼滤波\cite{anderson2001eakf}配合五阶分段有理局地化函数\cite{gaspari1999correlation}被广泛用于更新状态集合。然而现有研究尚未将集合同化样本集与双估计思想结合以解决元推理中的计算自评偏差。本文首次将同化更新的想象海洋划分为搜索用海洋与留出海洋，使卫星大数据标定的集合成员直接充当评判计算价值的外部标尺。
